\documentclass[10pt]{article}
\usepackage[margin=2cm]{geometry}
\usepackage{amsmath,amssymb,booktabs,graphicx,url,array}
\newcommand{\iseql}{ISEQL\textsuperscript{+}}
\newcommand{\tf}[1]{\texttt{#1}}
\newcommand{\rel}[1]{\mathit{#1}}
\newcommand{\op}[1]{\mathrm{#1}}
\newcommand{\lc}[1]{{}^{(#1)}\!}
\newcommand{\rc}[1]{^{(#1)}}
\setlength{\parskip}{3pt}\setlength{\parindent}{0pt}
\title{Supplementary Material\\[2pt]\large Toward Traceable Payment Fraud Features with ISEQL+ Interval Specifications}
\author{Trang Thuy Pham, Fabio Persia}
\date{Revision 6 (revision-4 evidence, repaired release, grouping sensitivity)}
\begin{document}
\sloppy
\maketitle

\noindent The scripts in \tf{analysis/} and \tf{engine/} generate every number below (see \tf{RUN.md} for the command order). They need the IEEE-CIS training files, obtained from the official Kaggle page and not redistributed, and the thesis pipeline in \tf{src/}, which produces \tf{train\_enriched.parquet}. Primary results use the revision-4 inputs in \tf{data4/}, with device preprocessing fitted on training data and frozen. Results labeled \emph{historical} come from earlier revisions and are not comparable. A clean-directory run of steps 2--7 of \tf{RUN.md} from \tf{train\_enriched.parquet} is logged in \tf{analysis/results/clean\_run.log}; every check script exits with a non-zero status on a primary mismatch, and \tf{MANIFEST.sha256} lists file hashes of the release. Feature-name aliases (thesis code names in parentheses in S1) differ only in labels.

\section*{S1. Parameter and provenance ledger}
\footnotesize
\begin{tabular}{@{}l p{3.4cm} p{6.6cm} p{4.0cm}@{}}
\toprule
ID & Source (type, locator) & Parameters (provenance) and conventions & Statistic (code alias) \\
\midrule
P1 & PCI SSC \& NCFTA bulletin, 21 Oct 2020 (industry guidance): low-value carding, p.~2; detection red flags, pp.~3--4; acquirer velocity checks, p.~5 & \$10 cutoff, 180\,s window, $p=0.6$, key \tf{card1} (all E). Rank proportion over distinct transaction IDs selected in wall-clock time; 0 if the history is empty. & \tf{micro\_tx\_ratio\_3min} \\
P2 & Engineering hypothesis, motivated by Reg.\ (EU) 2018/389 Art.\ 18(2)(c)(ii) (unusual device/software access; a condition for the risk-analysis exemption, not a prescribed session rule) & Specification: $p=0.3$ (E). Proxy: 30-day window (E). Device level: \tf{DeviceInfo} without the build suffix, lower-cased; levels with at least 20 \emph{training} occurrences are kept, all others (including unseen ones) become \tf{rare\_device}. The map is fitted once on training rows and frozen. A missing device is its own level (79.9\% of rows). The flag is 1 if the (group, level) pair has no strictly earlier observation within 30 days. & \tf{device\_novelty\_flag} \\
P3 & None (engineering hypothesis); the ACFE citation from the thesis was removed as unverified & Ratio 10, 3,600\,s, $p=0.8$ (E). Statistic: $\min(100,(\max-\min)/\mathrm{mean})$ of prior-hour amounts, 0 if fewer than 2 events. Amounts are finite and strictly positive (minimum \$0.251; asserted), so no zero-mean guard is triggered. The cap of 100 is an engineering choice and is never reached (maximum 13.8). & \tf{amount\_spread\_1h} (\tf{escalation\_ratio\_1h}) \\
P4 & 31 CFR 1010.311 and 1010.313; 31 U.S.C.\ 5324 (cash reporting and structuring; analogy only) & \$9,500 cutoff (E, inspired by the \$10,000 currency threshold), $k=3$ (E). Buckets are 24-hour elapsed-time buckets starting at $t_0=86{,}400$\,s (the first \tf{TransactionDT}), not business days. Primary statistic: strictly earlier events in the anchor's bucket. & \tf{subthreshold\_count\_elapsed} (\tf{structuring\_count\_elapsed}); batch: \tf{structuring\_count\_day} \\
P5 & PCI SSC \& NCFTA bulletin, p.~5 (velocity checks) & Windows of 1, 6 and 24\,h; $k=5$ (E). & \tf{tx\_count\_\{1h,6h,24h\}} \\
P6 & PCI DSS v4.0 Req.\ 8.3.4 (system-access lockout after at most 10 invalid attempts); Reg.\ 2018/389 Art.\ 4(3)(b) (at most 5 consecutive failed authentications) & $k=10$ (S, transferred by analogy); $\delta=300$\,s (E); proxy cutoff \$1 (E). The transaction proxy does not implement either control. & \tf{small\_amount\_count\_5min} (\tf{micro\_fail\_count\_5min}) \\
\bottomrule
\end{tabular}
\normalsize

\section*{S2. Preprocessing, conformance and prefix invariance}
\textbf{Device preprocessing (revision 4).} The map is fitted on the training split and frozen. Compared with the earlier whole-log fit, 1,687 device values change (934 / 256 / 497 in train / validation / test) and 798 P2 flags change (499 / 126 / 173). A synthetic fixture (S4) shows why a whole-log fit can alter an earlier novelty flag when future rows are appended, while a frozen training map cannot. Category dictionaries for LightGBM cover the whole log. These are vocabulary metadata, not fitted statistics: levels absent from training receive no splits of their own. We therefore claim only that the device-level map (and hence P2) is fitted on training data, not that every part of preprocessing is.

\textbf{Reference comparison} (\tf{conformance\_v4.py}). This is an independently written reference implementation of the declared \emph{observable} statistics (novelty, capped spread and small-payment count for P2, P3, P6; it does not execute the session, escalation or authentication queries): a brute-force, per-row loop over each group's events. The sample is all 62 \tf{card1} groups containing tied timestamps plus 300 random groups (seed 7): 362 groups, 178,237 transactions, of which 312 are tied. The script asserts unique IDs, one-to-one validated joins, an unchanged row count, and the amount domain. Exclusions are derived (rows without a finite implemented value), not hard-coded, and the script fails unless all primary statistics have 0 mismatches and 0 exclusions.

\begin{center}\footnotesize
\begin{tabular}{@{}lrrrrrl@{}}
\toprule
Statistic & Tested & Empty/default cases & Matches & Mismatches & Excluded & Comparison \\
\midrule
P1 & 178,237 & 150,560 empty & 178,237 & 0 & 0 & abs.\ tol.\ $10^{-9}$, rel.\ 0 \\
P2 (frozen map, tie-corrected) & 178,237 & -- & 178,237 & 0 & 0 & exact \\
P3 & 178,237 & 85,116 with $<$2 events & 178,237 & 0 & 0 & abs.\ tol.\ $10^{-9}$ \\
P4 elapsed / P4 batch & 178,237 & 12,926 empty / -- & 178,237 & 0 & 0 & exact \\
P5 (1\,h, 6\,h, 24\,h), P6 & 178,237 each & -- & 178,237 & 0 & 0 & exact \\
\bottomrule
\end{tabular}\end{center}

\textbf{P2 tie history.} In revision 2, the case was \tf{card1} 2380, IDs 3322544 and 3322545, at the same \tf{TransactionDT} and both without device. A \tf{shift} on rows used the tied peer as history. All P2 values since revision 3 shift on distinct timestamps.

\textbf{End-to-end prefix check} (\tf{prefix\_v4.py}). The script starts from the raw fields (\tf{TransactionDT}, \tf{card1}, \tf{TransactionAmt}, \tf{DeviceInfo}) and the frozen map. Recomputing all statistics on the full log reproduces the saved revision-4 vectors with 0 mismatches for each statistic. Recomputing on logs truncated at seven recorded cut times gives the counts of changed earlier values below.

\begin{center}\footnotesize
\begin{tabular}{@{}llrrcc@{}}
\toprule
Cut & Mode & \tf{TransactionDT} & Rows kept & P1--P6 (all) & Batch P4 \\
\midrule
start of validation (day 120) & $t<$ cut & 10,454,400 & 414,542 & 0 & 0 \\
start of test (day 150) & $t<$ cut & 13,046,400 & 498,113 & 0 & 0 \\
mid-day (day 160 + 12,345\,s) & $t<$ cut & 13,922,745 & 529,396 & 0 & 323 \\
first tied timestamp in test & $t<$ cut & 13,207,022 & 504,327 & 0 & 1,813 \\
tied timestamp \#1, batch included & $t\le$ cut & 13,207,022 & 504,329 & 0 & 1,815 \\
tied timestamp \#2, batch included & $t\le$ cut & 13,711,497 & 521,729 & 0 & 1,040 \\
tied timestamp \#3, batch included & $t\le$ cut & 14,061,910 & 533,654 & 0 & 1,169 \\
\bottomrule
\end{tabular}\end{center}
These checks verify post-training scoring-time availability under the frozen training map; the script fails if any primary statistic changes. They do not claim that each \emph{training} feature was constructible at its own timestamp, because the map uses the whole training period.

\section*{S3. Set-level coverage}
\textbf{Definition.} $\nu_0=R\,\omega_g\,S$ is the unconstrained join with group equality, and $\nu_0(r)=\{s:\langle r,s\rangle\in\nu_0\}$; the group restriction precedes union formation. Then $C(r)=\#\big(r\cap\bigcup_{s\in\nu_0(r)}s\big)/|r|$ on integer chronons, for a nonempty closed anchor, and $R\,{}^{C(\theta''p)}\omega_g\,S=\{\langle r,s\rangle\in\nu_0 : C(r)\,\theta''\,p\}$, $p\in[0,1]$. Other constraints decorating the same operator refer to the same $\nu_0$. It returns no pair for an anchor without partners. The scalar feature is total (0 in that case), so the alert equivalence is stated only for positive lower-bound thresholds. An anchor-preserving aggregate output would need a separate operator.

\textbf{Computation} (\tf{coverage\_tests.py}). Restrict to the group, clip each partner to $r$, sort, coalesce overlapping or adjacent intervals, and sum their lengths; this takes $O(m\log m)$. Union coverage is an elementary set measure; the proposal concerns its use as a grouped \iseql{} constraint.

\textbf{Properties and proof.} Let $K$ be the chronons of $r$, $A_s=K\cap s$ and $U=\bigcup_s A_s$, so $C(r)=\#U/\#K$.
\begin{itemize}
\item $0\le C\le1$, since $U\subseteq K$ and $K\neq\emptyset$.
\item Repeating a partner leaves $U$ unchanged; adding one enlarges or preserves $U$, so $C$ never decreases (an upper-bound predicate need not be preserved).
\item Let $\kappa(x)=\#\{s: x\in A_s\}$. Then $\sum_s\#A_s-\#U=\sum_{x\in U}(\kappa(x)-1)\ge 0$, with equality if and only if no chronon lies in two clipped partners, i.e., the nonempty clipped partners are pairwise disjoint (both directions).
\end{itemize}
These are elementary; they do not establish novelty of union coverage, only its behavior as a grouped \iseql{} constraint.
The hand cases are:

\begin{center}\footnotesize
\begin{tabular}{@{}llccc@{}}
\toprule
Partners of $r=[1,10]$ & Case & $C(r)$ & Normalized sum & Sum capped at 1 \\
\midrule
{}[1,8], [3,10] & overlapping & 1.00 & 1.60 & 1.00 \\
{}[2,9], [4,5] & nested & 0.80 & 1.00 & 1.00 \\
{}[1,3], [6,7] & disjoint & 0.50 & 0.50 & 0.50 \\
{}[2,5], [2,5] & duplicate copies & 0.40 & 0.80 & 0.80 \\
{}[4,4], [4,4] & coincident points & 0.10 & 0.20 & 0.20 \\
-- & empty & no pair (feature 0) & 0 & 0 \\
{}[8,15] & beyond anchor & 0.30 & 0.30 & 0.30 \\
\bottomrule
\end{tabular}\end{center}
The nested case shows that capping the sum at 1 does not repair it as a coverage measure. Group restriction matters: a cross-group partner changes $C$ from 0.2 to 0.5 if the grouping filter is skipped. A randomized check of 20,000 anchors confirms the union definition, bounds, duplicate invariance, monotonicity and both directions of the equality condition.

\textbf{Proposition (pairwise ratios do not determine $C$).} For a pair $\langle r,s\rangle$ write the left ratio $\lambda(r,s)=\#(r\cap s)/|r|$ and the right ratio $\varrho(r,s)=\#(r\cap s)/|s|$. Define the \emph{pairwise profile} of an instance as the tuple $\big(r,\ |\nu_0(r)|,\ \{\!\{\lambda(r,s)\}\!\},\ \{\!\{\varrho(r,s)\}\!\}\big)$, the anchor, the partner count and the two multisets of pairwise ratios. Then $C(r)$ is not a function of the pairwise profile. Consequently no constraint computed from pairwise overlap ratios on the same operator -- under any boolean combination or aggregate, including the sum, the sum capped at 1, extrema and order statistics -- expresses $C(r)\,\theta''\,p$.

\emph{Proof.} It suffices to exhibit two instances with equal pairwise profiles and different $C$. Any function of the profile takes the same value on both, while $C$ does not. The witnesses below use partners that are pairwise distinct and of equal length, so the construction does not rely on duplicates or degenerate intervals. \hfill$\square$

\begin{center}\footnotesize
\begin{tabular}{@{}llcccc@{}}
\toprule
$r$ & Partners & $\{\!\{\lambda\}\!\}$ & $\{\!\{\varrho\}\!\}$ & $C(r)$ & Norm.\ sum (capped) \\
\midrule
{}[1,10] & [1,5], [6,10]            & $\{1/2,1/2\}$ & $\{1,1\}$ & 1.00 & 1.00 (1.00) \\
{}[1,10] & [1,5], [2,6]             & $\{1/2,1/2\}$ & $\{1,1\}$ & 0.60 & 1.00 (1.00) \\
{}[1,12] & [1,4], [5,8], [9,12]     & $\{1/3,1/3,1/3\}$ & $\{1,1,1\}$ & 1.00 & 1.00 (1.00) \\
{}[1,12] & [1,4], [2,5], [3,6]      & $\{1/3,1/3,1/3\}$ & $\{1,1,1\}$ & 0.50 & 1.00 (1.00) \\
\bottomrule
\end{tabular}\end{center}

\textbf{Corollary (the discrepancy can approach one).} Let $|r|=k^2$. Tiling $r$ with the $k$ blocks $[(i-1)k+1,\,ik]$ and crowding it with the $k$ unit-shifted intervals $[i,\,i+k-1]$, $i=1,\dots,k$, give equal pairwise profiles with $C=1$ and $C=(2k-1)/k^2$. The discrepancy is $1-(2k-1)/k^2\to 1$; since both coverages lie in $[0,1]$ it is bounded by 1, and the construction shows it can be made arbitrarily close to that maximum rather than unbounded. Verified for $k=2,\dots,8$: gaps 0.250, 0.444, 0.563, 0.640, 0.694, 0.735, 0.766.

The proposition is a non-equivalence result about pairwise overlap-ratio constraints on one operator. It is \emph{not} a claim that SQL, a general \iseql{} query composition, or any other formalism cannot compute union coverage; ordinary window aggregates can. All witnesses, properties and the corollary are checked in exact rational arithmetic by \tf{engine/test\_inexpressibility.py} and \tf{engine/test\_ops.py}, which exit non-zero on failure.

\textbf{Rank reading of P1 (anchor-specific).} For each anchor $(g,t)$, first select $H_{180}(t)$ in wall-clock time; if $n>0$, let $\rho_{(g,t)}$ be a bijection from the distinct IDs in $H_{180}(t)$ to $\{1,\dots,n\}$. Replace the anchor by $[1,n]$ and each event $j$ by $[\rho(j),\rho(j)]$. Filtering micro-transactions then gives $C=m/n$. Wall-clock selection and the multiplicity of distinct events are preserved; durations, gaps and wall-clock occupancy are not. Two distinct IDs at the same second count twice on ranks but occupy one wall-clock chronon (fixture in S4). P1 is exact \emph{under this declared reading}; this is not a general theorem that an event clock preserves interval joins. Input contract: IDs are unique (asserted); a duplicated ID is rejected, not de-duplicated.

\section*{S4. Specification fixtures} (\tf{spec\_fixtures\_v4.py}; all pass; exits non-zero on failure)
\begin{itemize}\itemsep1pt
\item \textbf{P3 escalation spans} are typed tuples $\rel{Esc}(id_i,id_k,g,t_i,t_k)$ for same-group ordered pairs with $t_i<t_k$ and $a_k>10a_i$. They involve no segmentation and no monotonicity of intermediate amounts. A span qualifies for anchor $t$ if it ends inside $w=[t-3600,t-1]$ (LOJ $\cup$ DJ) and at least 80\% of its chronons lie inside $w$.
\begin{itemize}
\item Amounts [1, 2, 20] qualify; the reversed [20, 2, 1] do not; both have spread 2.478.
\item A cross-group pair is not a span.
\item A span with about 67\% of its own chronons inside $w$ is excluded by the 80\% left-overlap constraint.
\item The fixture clock is shifted so that all window endpoints are nonnegative.
\end{itemize}
\item \textbf{P6 timestamp-batch reset.} For a success at $t_s$, the eligible failures satisfy $t_s-\delta\le t_f<t_s$ and $t_f>$ the latest earlier success time. With F@30, S@30, F@40, S@70, only F@40 counts. In F,F,S,F,F,F,S@70 plus a tied F@70, the eligible failures are \{40, 50, 60\}. IEEE-CIS has no authentication log, so P6 is never executed on data.
\item \textbf{P1 rank axis.} Two tied micro events plus one large event give $m/n=2/3$ on ranks but one wall-clock chronon; an input with a duplicated ID is rejected.
\item \textbf{Empty partners:} no pair is returned, even though the scalar default 0 satisfies an upper bound.
\item \textbf{Device pooling.} Using the review's counterexample, fitting the map on the whole log flips the earlier novelty flag from 0 to 1 once 19 future rows are appended; the frozen earlier-period map does not.
\end{itemize}

\section*{S5. Replayable trace fixtures} (\tf{traces/*.json}; \tf{python3 replay\_trace.py traces/*.json})
Each fixture contains:
\begin{itemize}\itemsep0pt
\item the source and locator, specification version, parameters with their provenance, and conventions;
\item the anchor;
\item every nearby same-group event with its relative time, amount and eligibility reason;
\item the stored feature value;
\item the label and D score, kept separate from the rule alert.
\end{itemize}
Not recorded in the fixtures: the data-snapshot and preprocessing hashes (listed in \tf{MANIFEST.sha256}) and a model operating threshold (no model decision is claimed). The replayer verifies the witness from the fixture's events; checking that no eligible event was omitted requires the source snapshot. It exits non-zero on any mismatch. Replay output:
\begin{center}\footnotesize
\begin{tabular}{@{}lccccl@{}}
\toprule
Fixture & $n$ & $m$ & $\varphi_1$ (replayed = stored) & Alert & Label / D score \\
\midrule
alert, single prior event (3575098) & 1 & 1 & 1.000 & yes & fraud / 0.901 \\
boundary event at exactly $-180$\,s, included (3537612) & 1 & 0 & 0.000 & no & legitimate / 0.026 \\
non-alert at 0.5 (3494766) & 2 & 1 & 0.500 & no & legitimate / 0.271 \\
\bottomrule
\end{tabular}\end{center}
A synthetic illustration (not data): a group-A anchor at $t=200$ with \$150 has earlier events \$4 ($t=20$), \$7 ($t=100$) and \$30 ($t=180$). A group-B event at 150 is ineligible, a group-A event at 200 is tied and excluded, and one at 210 is unavailable. So $\varphi_1=2/3$ and the rule fires although the current amount is large.

\section*{S6. Configuration}
\begin{itemize}\itemsep1pt
\item \textbf{Split} by elapsed day from $t_0$: training $<120$ (414,542 rows), validation 120--149 (83,571), test $\ge150$ (92,427).
\item \textbf{Features:} 436 standard columns (listed in \tf{models/column\_meta.json}): \tf{TransactionAmt}, \tf{dist1/2}, screen size, C1--C14, D1--D15, V1--V339, \tf{id\_01}--\tf{id\_11}, 34 categoricals (including \tf{card1} and the device level), 17 missingness indicators, \tf{has\_identity}. Model D adds the 8 statistics; D\textsubscript{b} replaces P4 by its batch count. Numeric missing values are set to 0, and LightGBM handles categoricals natively.
\item \textbf{LightGBM} 4.7.0: \tf{num\_leaves}=31, \tf{learning\_rate}=0.05, \tf{n\_estimators}=500, \tf{is\_unbalance}, \tf{metric}=AUC with \tf{first\_metric\_only}, early stopping with patience 50, seed 42, no bagging or feature subsampling.
\begin{itemize}
\item Best iterations: C 228, D 215, D\textsubscript{b} 264.
\item Ablations: $-$P1 215, $-$P2 258, $-$P3 215, $-$P4 294, $-$P5 198, $-$P6 215.
\item Training is deterministic: refits without P1 or P6 reproduce D's predictions exactly.
\end{itemize}
\item \textbf{Deletion probe:} the removed legitimate training IDs are 3182001, 3333738, 3322921 and 3287655 (\tf{data4/drop\_manifest.json}, seeds 1--4). Histories and preprocessing are held fixed; only the training row is omitted.
\item \textbf{Bootstrap:} paired percentile bootstrap over \tf{card1} groups (6,182 in validation, 6,402 in test). Each replicate draws groups with replacement, includes all rows of each drawn group, and recomputes AUC and average precision on the same draw for both models; 1,000 replicates, seed 42. No replicate lacked either class. Intervals condition on the fitted models and on the partition, not on training variability; we treat them as primary for these grouped data. DeLong tests assume independent transactions.
\item \textbf{SHAP} 0.52.0 TreeExplainer: path-dependent, no background set, log-odds output, additivity error $4.4\times10^{-14}$. The global ranking uses 1,000 validation rows (seed 42). The score-binned shares use all 3,213 test frauds plus 6,000 random legitimate rows, with labels joined by \tf{TransactionID} (the CSV carries \tf{TransactionID} and \tf{isFraud}); the script asserts that its feature list equals the stored model's names and order; $q_i=\sum_{\text{stat}}|\phi|/\sum_{\text{all}}|\phi|$, and no zero denominators occurred.
\item \textbf{Software:} Python 3.13.16, LightGBM 4.7.0, pandas 3.0.5, NumPy 2.5.3, scikit-learn 1.9.1, SciPy 1.18.1, PyArrow 25.0.1, SHAP 0.52.0, Matplotlib 3.11.2 (\tf{requirements.txt}); 2 CPU threads (LightGBM default).
\item \textbf{Historical (not current) LR results:} revision-1 models A and B with original features gave validation AUC 0.8502 / 0.8499 and test 0.8301 / 0.8296. They were not refitted with the corrected statistics and are not part of the current design.
\end{itemize}

\section*{S7. Additional results (revision 4)}
\textbf{Ablation} (contribution $=$ AUC(D) $-$ AUC(D without pattern), $\times10^{-3}$, with paired 95\% group intervals). These are exploratory single fits.
\begin{center}\footnotesize
\begin{tabular}{@{}lcc@{}}
\toprule
Pattern & Validation & Test \\
\midrule
P2 & $-0.84$ [$-2.05$, $+0.29$] & $+1.90$ [$+0.49$, $+3.38$] \\
P3 & $-0.69$ [$-1.66$, $+0.21$] & $+0.09$ [$-0.81$, $+1.01$] \\
P4 & $-0.68$ [$-2.20$, $+0.88$] & $+2.24$ [$+0.67$, $+3.91$] \\
P5 & $-0.80$ [$-2.64$, $+0.72$] & $+0.06$ [$-2.27$, $+2.46$] \\
P1, P6 & 0 (refit identical) & 0 (refit identical) \\
\bottomrule
\end{tabular}\end{center}

\begin{center}\includegraphics[width=0.5\textwidth]{figures/figS1_ablation.pdf}\\\footnotesize Fig.\ S1. Exploratory ablation with paired group intervals (* proxy).\end{center}

\textbf{Deletion probe} (D $-$ C, AUC).
\begin{center}\footnotesize
\begin{tabular}{@{}lcccc@{}}
\toprule
Fit & C val & D val & D$-$C val & D$-$C test \\
\midrule
original & 0.9145 & 0.9133 & $-0.0012$ & $-0.0005$ \\
drop 1 & 0.9152 & 0.9149 & $-0.0003$ & $-0.0015$ \\
drop 2 & 0.9157 & 0.9141 & $-0.0016$ & $-0.0041$ \\
drop 3 & 0.9154 & 0.9150 & $-0.0004$ & $-0.0025$ \\
drop 4 & 0.9156 & 0.9142 & $-0.0014$ & $-0.0001$ \\
\bottomrule
\end{tabular}\end{center}

\textbf{P2 strata} (test month; lift with 95\% group interval).
\begin{center}\footnotesize
\begin{tabular}{@{}lrrrrl@{}}
\toprule
Stratum & Rows & Fraud rate & Alerts & TP & Lift [CI] \\
\midrule
observed device, established group & 15,428 & 9.2\% & 3,376 & 298 & 0.96 [0.78, 1.19] \\
observed device, first observation of group & 235 & 10.2\% & 235 & 24 & 1 (all flagged) \\
missing device, established group & 76,350 & 2.3\% & 1,650 & 30 & 0.79 [0.53, 1.10] \\
missing device, first observation of group & 414 & 1.9\% & 414 & 8 & 1 (all flagged) \\
\bottomrule
\end{tabular}\end{center}
The overall lift is 1.82 [1.58, 2.10] (1.8248). Within the two established-group strata, novelty alerts do not demonstrate enrichment over their stratum base rates; the first-observation strata are entirely flagged by definition. \emph{Composition benchmark} (\tf{diagnostics\_v4.py}): the expected number of true positives if alerts had their stratum's fraud rate is $\sum_k a_k r_k=381.8$ against 360 observed, ratio 0.94 [0.79, 1.13] (same 1,000 \tf{card1}-group draws). The marginal enrichment is thus accounted for descriptively by device-availability composition; this is an association, not a causal decomposition.

\textbf{Section V diagnostics} (\tf{diagnostics\_v4.py}; current revision-4 features only).
\begin{itemize}\itemsep0pt
\item P1, test month: 216 alerts, 22 TP; eligible history $n=1$: 152 alerts, 22 TP; $n\ge2$: 64 alerts, 0 TP.
\item Class means, full log (590,540 rows, unweighted): five of eight statistics are lower for fraud (P3, P4, three P5 counts); P1, P2 and P6 are higher. Example: \tf{tx\_count\_24h} 15.28 vs.\ 18.93.
\item Quintiles of transactions ranked by whole-log \tf{card1} group size (ties broken by chronological order; boundaries 79, 476, 1,538, 3,977 transactions; retrospective description only): mean \tf{tx\_count\_24h} fraud vs.\ legitimate 0.89/0.42, 2.20/1.67, 8.28/10.23, 15.59/16.68, 46.89/65.80.
\item Alert rates (test): P4 45.0\%, P4\textsuperscript{b} 69.2\%, P5 6.8/28.2/51.4\% (1/6/24\,h), P2 6.1\%, P1 0.2\%. P6 has no alerts, so its precision and lift are not estimable; this does not show that the behavior is absent.
\end{itemize}

\textbf{SHAP (revision-4 D).}
\begin{itemize}
\item Global ranks among 444 features: \tf{card1} 1, \tf{card2} 2, \tf{V258} 19; \tf{tx\_count\_24h} 32, \tf{tx\_count\_1h} 67, \tf{device\_novelty\_flag} 74, \tf{tx\_count\_6h} 121, P3 131, P4 180. P1 and P6 have zero attribution.
\item A statistic is among the top five drivers for 10 of 3,213 test frauds, all \tf{tx\_count\_24h}.
\item Mean $q_i$ is 1.07\% for frauds scored below 0.3 and 0.48\% for those scored at least 0.9 (Spearman $\rho=-0.52$). Mean signed contribution: $+0.019$ (fraud) and $+0.010$ (legitimate) log-odds.
\end{itemize}

\begin{center}\includegraphics[width=0.48\textwidth]{figures/figS2_attribution.pdf}\\
\footnotesize Fig.\ S2. Mean $q_i$ (\%) for test frauds by D score; the values are descriptive, with no interval shown (white numbers are cases per bin). Scores are uncalibrated outputs of a class-reweighted classifier and should not be interpreted as calibrated real-world fraud probabilities.\end{center}

\clearpage
\section*{S8. Grouping sensitivity (one composite partition)}
\textbf{Protocol.} Written and hashed after the original \tf{card1} diagnostics motivated a sensitivity check, but before any composite-key statistic, rule evaluation or model result was computed (\tf{analysis/results/PROTOCOL\_uid.md}, SHA-256 \tf{0283506\ldots}). Key: $\mathit{uid}=(\tf{card1},\tf{addr1},\text{elapsed day}-\tf{D1})$ with raw \tf{D1} (no imputation); a missing \tf{addr1} or \tf{D1} is its own level of that component. Only this partition was tried. It is an engineering partition, not a recovered cardholder. Same eight specifications, conventions and defaults; P2 uses the frozen device map. Primary comparison: D\textsubscript{uid} vs.\ C on the test month, scored once; claim rule: test $\Delta$AUC $>0$ with a group interval excluding 0 and a positive sign in all five probe fits. The rule-level and direction summaries below were added afterwards, for the same single partition.

\textbf{Checks.} With \tf{card1} as key, \tf{stats\_by\_key.py} reproduces all eight revision-4 vectors on all 590,540 rows (0 mismatches), and \tf{grouping\_sensitivity.py} reproduces Table III. The reference implementation on $\mathit{uid}$ groups (all 51 groups with tied timestamps plus 300 random groups; 1,783 rows, 132 tied) gives 0 mismatches for all eight statistics. Bootstrap unit: \tf{card1} for both partitions ($\mathit{uid}$ groups are nested in \tf{card1}).

\begin{center}\footnotesize
\begin{tabular}{@{}lrr@{}}
\toprule
 & \tf{card1} & $\mathit{uid}$ \\
\midrule
groups & 13,553 & 217,850 \\
median group size (per group / per transaction) & 4 / 919 & 1 / 5 \\
largest group & 14,932 & 1,414 \\
rows with no earlier same-group event in 24\,h & 21.6\% & 71.0\% \\
statistics lower for fraud (class means, full log) & 5 of 8 & 1 of 8 (P2) \\
mean \tf{tx\_count\_24h}, fraud vs.\ legitimate & 15.28 vs.\ 18.93 & 2.34 vs.\ 1.96 \\
mean \tf{subthreshold\_count\_elapsed} & 8.06 vs.\ 10.11 & 2.04 vs.\ 1.56 \\
mean \tf{amount\_spread\_1h} & 0.269 vs.\ 0.293 & 0.053 vs.\ 0.027 \\
\bottomrule
\end{tabular}\end{center}

\textbf{Rules, test month} (lift with 95\% \tf{card1}-group interval).
\begin{center}\footnotesize
\begin{tabular}{@{}lrrlrrl@{}}
\toprule
 & \multicolumn{3}{c}{\tf{card1}} & \multicolumn{3}{c}{$\mathit{uid}$} \\
\cmidrule(lr){2-4}\cmidrule(l){5-7}
Rule & Alerts & TP & Lift & Alerts & TP & Lift \\
\midrule
P1 & 216 & 22 & 2.93 [0.92, 7.26] & 185 & 16 & 2.49 [0.63, 6.73] \\
P2* & 5,675 & 360 & 1.82 [1.58, 2.10] & 39,842 & 1,326 & 0.96 [0.88, 1.03] \\
P4 & 41,590 & 1,501 & 1.04 [0.93, 1.15] & 8,092 & 512 & 1.82 [1.29, 2.53] \\
P5 1\,h & 6,275 & 152 & 0.70 [0.39, 1.27] & 1,698 & 20 & 0.34 [0.06, 2.56] \\
P5 6\,h & 26,064 & 907 & 1.00 [0.81, 1.22] & 3,191 & 166 & 1.50 [0.62, 3.61] \\
P5 24\,h & 47,506 & 1,707 & 1.03 [0.92, 1.15] & 5,209 & 346 & 1.91 [1.13, 3.27] \\
P6* & 0 & 0 & -- & 0 & 0 & -- \\
\bottomrule
\end{tabular}\end{center}
Under $\mathit{uid}$, P2's first-observation default fires on most singleton groups (39,842 alerts), and within the established-group strata its lift is 0.96 [0.73, 1.35] (observed device) and 0.54 [0.41, 0.67] (missing device).

\textbf{Model} (pre-registered; \tf{train\_uid.py}, \tf{eval\_uid.py}). D\textsubscript{uid} = C + the eight $\mathit{uid}$ statistics, identical configuration.
\begin{center}\footnotesize
\begin{tabular}{@{}lcc@{}}
\toprule
 & Validation & Test \\
\midrule
$\Delta$AUC vs.\ C, DeLong $p$ & $-0.0006$, $p=0.10$ & $-0.0012$, $p<0.001$ \\
\quad 95\% group CI & [$-0.0020$, $0.0009$] & [$-0.0027$, $0.0003$] \\
\quad original + 4 deletion refits & [$-0.0007$, $+0.0002$] & [$-0.0033$, $+0.0005$] \\
$\Delta$PR-AUC vs.\ C [95\% group CI] & $+0.0049$ [$0.0014$, $0.0086$] & $+0.0012$ [$-0.0025$, $0.0057$] \\
$\Delta$AUC vs.\ D (\tf{card1} statistics) & $+0.0006$, $p=0.24$ & $-0.0006$, $p=0.31$ \\
\bottomrule
\end{tabular}\end{center}
The claim rule is not met: no predictive gain is demonstrated for either partition. The counter-intuitive directions and the null rule results for P4 and the 24-hour P5 are therefore partition-specific, whereas the absence of a model gain is not. Neither outcome identifies a mechanism or verifies identities.
\clearpage
\section*{S9. Interval-join evaluation of the exactly expressible specifications}
(\tf{engine/}; run order in \tf{engine/RESULTS\_engine.md})

\textbf{What is executed.} The operators are implemented directly: the during join $\op{DJ}$ with a cardinality constraint, and its reverse $\op{RDJ}$ with the set-level coverage constraint $C$ of S3, with group equality folded into the join predicate. Two independent code paths evaluate the same operators. \tf{reference\_eval*} materializes every window interval and every partner interval and computes $C$ by clipping partners to the anchor, sorting, coalescing and summing chronons. \tf{sweep\_*} sorts the relation once by (group, start point) and sweeps each group with monotone endpoints, one pass per operator. Direct feature-vector conformance is evaluated for P1, P4 and P5, whose stored statistics are exact under the declared rank or elapsed-time conventions. A separate semantic experiment below also executes P3's conceptual escalation-span predicate; it is compared with, but is not represented by, the lossy \tf{escalation\_ratio\_1h} statistic, so no feature vector is reproduced for it. P2 and P6 remain unexecuted conceptual predicates, because the data records neither the sessions nor the authentication events they quantify over.

\textbf{Agreement.} The sweep reproduces the manually derived vectors on all 590{,}540 rows, and the materializing implementation agrees with both on the S2 sample (362 groups, 178{,}237 transactions, all 62 tied-timestamp groups).

\begin{center}\footnotesize
\begin{tabular}{@{}llcc@{}}
\toprule
Statistic & Specification evaluated & Rows & Disagreement \\
\midrule
\tf{micro\_tx\_ratio\_3min}        & $w_{180}\ {}^{C_{\mathrm{rank}}(\geq 0.6)}\!\op{RDJ}\ \sigma_{\rel{amt}<10}(\rel{Tx})$ & 590,540 & max abs.\ diff.\ 0 \\
\tf{subthreshold\_count\_elapsed}  & $\sigma_{\rel{amt}<9500}(\rel{Tx})\ \op{DJ}\rc{\geq 3}\ [d(t),t{-}1]$ & 590,540 & 0 rows \\
\tf{tx\_count\_1h}                 & $\rel{Tx}\ \op{DJ}\rc{\geq 5}\ w_{3600}$  & 590,540 & 0 rows \\
\tf{tx\_count\_6h}                 & $\rel{Tx}\ \op{DJ}\rc{\geq 5}\ w_{21600}$ & 590,540 & 0 rows \\
\tf{tx\_count\_24h}                & $\rel{Tx}\ \op{DJ}\rc{\geq 5}\ w_{86400}$ & 590,540 & 0 rows \\
\bottomrule
\end{tabular}\end{center}

\textbf{Scope of the claim.} This is agreement between two algorithms evaluating one declared set of conventions, not a second independent derivation in the sense of S2, whose reference implementation was written independently of the original. It establishes that the specifications as written compute the stored vectors; it does not validate the conventions themselves, and it is not a production engine.

\textbf{Cost (illustrative, not a benchmark).} On one 4-core x86-64 Linux container with Python 3.10, \tf{numpy} 2.2.6 and \tf{pandas} 2.3.3, a single warm run of all five specifications over 590,540 anchors took 1.8\,s in total (sort 0.06\,s; P1 0.59\,s; P4 0.27\,s; P5 0.28--0.29\,s per window), against 4.4\,s for the \tf{pandas} derivation. Both figures exclude parquet I/O and preprocessing, are single unrepeated runs on one machine, and are reported to show feasibility rather than as a performance comparison. Per anchor the sweep uses monotone endpoints, and coverage costs $O(m\log m)$ in the partner count (S3).

\textbf{Additional semantic checks.} Three questions the main text can only state are evaluated here.

\emph{The burst composition.} $w_{180}\ {}^{(\geq 3),\,C_{\mathrm{rank}}(\geq 0.6)}\!\op{RDJ}\ \sigma_{\rel{amt}<10}(\rel{Tx})$ requires dominance \emph{and} repetition. On the test month, requiring at least two micro-partners leaves 64 alerts and \emph{no} true positive; at least three leaves 9 alerts and none; at least five leaves none. Of P1's 216 alerts, 152 have a history of exactly one event, and all 22 true positives lie there. The faithful burst reading therefore removes every true positive on this data, which is why the paper reports P1 as a dominance test rather than a weaker burst detector.

\emph{P1 on wall-clock time.} With point-interval partners $C_{\mathrm{time}}(w)\le n/|w|$, so $C_{\mathrm{time}}\ge 0.6$ over 180 chronons needs 108 covered seconds. The largest 180\,s history in the data has $n=21$ (ceiling 0.117); explicit evaluation of all 51,996 anchors with a non-empty history gives a maximum of 0.022 and no alert at any threshold at or above 0.05. The rank bijection of S3 is therefore load-bearing: without it the constraint is unsatisfiable at the declared $p$.

\emph{P3's faithful span specification.} $\rel{Esc}\ \lc{\geq 0.8}(\op{LOJ}\cup\op{DJ})\ w_{3600}$ can be evaluated even though its witness statistic is a lossy summary (candidate spans satisfy $\mathrm{inside}\le 3600$ and $\mathrm{inside}/L\ge 0.8$, hence $L\le 4500$, so a 4{,}500\,s lookback is exact). On the test month it yields 2,445 alerts and 68 true positives, precision 2.78\%, lift 0.80; allowing the anchor to close the span yields 3,676 alerts, 120 true positives, lift 0.94, and the two readings differ on 1,231 anchors. Both are below the base rate. At a count-matched threshold the specification and \tf{escalation\_ratio\_1h} share only 1,587 alerts (Jaccard 0.48, 858 specification-only and 859 statistic-only); 15,652 anchors have a non-zero statistic with no qualifying span and 63 the converse. The summary is therefore not a conservative approximation of the specification in either direction.

\end{document}
